\begin{examplebox}

\textbf{\textcolor{CExample}{例 1（基础：已知横坐标，直接代入）}}

\textbf{题目：}
已知双曲线 $\frac{x^2}{16}-\frac{y^2}{9}=1$，点
$P\left(5,\frac94\right)$ 在双曲线上，求 $|PF_1|$ 和 $|PF_2|$。

\textbf{解：} $a=4$、$b=3$、$c=5$，故 $e=\frac54$。
先验证点在曲线上：
\[
 \frac{5^2}{16}-\frac{(9/4)^2}{9}
 =\frac{25}{16}-\frac9{16}=1.
\]
$x_0=5\ge4$，在右支，所以
\[
 \begin{aligned}
 |PF_1|&=\frac54\cdot5+4=\boxed{\frac{41}{4}},\\
 |PF_2|&=\frac54\cdot5-4=\boxed{\frac94}.
 \end{aligned}
\]
\textbf{自检：} $\frac{41}{4}-\frac94=8=2a$。

\medskip
\textbf{\textcolor{CExample}{例 2（逆向求方程：原题保留并补足检验）}}

\textbf{题目：}
焦点在 $x$ 轴的双曲线右支上一点 $P(5,y_0)$，
到\textbf{左、右焦点}的距离分别是 $13$ 和 $3$，求双曲线方程。

\textbf{解：} 右支长短分别为 $ex_0+a$ 与 $ex_0-a$，所以
\[
 \begin{cases}5e+a=13,\\5e-a=3.\end{cases}
 \quad\Longrightarrow\quad
 10e=16,\ \ 2a=10.
\]
故 $e=\frac85$，$a=5$，$c=ae=8$，
$b^2=c^2-a^2=64-25=39$。
\[
 \boxed{\frac{x^2}{25}-\frac{y^2}{39}=1}.
\]
\textbf{易忽略的验证：}
本题 $x_0=5=a$，所以 $P$ 恰在\textbf{右顶点}，$y_0=0$。
实际距离为 $13=8+5$、$3=8-5$，全部吻合。

\medskip
\textbf{\textcolor{CExample}{例 3（焦半径之比：先判断所在支）}}

\textbf{题目：}
$F_1,F_2$ 是双曲线 $\frac{x^2}{16}-\frac{y^2}{9}=1$ 的左右焦点。
点 $P$ 在曲线上且 $|PF_1|:|PF_2|=3:1$，求 $P$ 的坐标。

\textbf{解：} $a=4$、$c=5$、$e=\frac54$。
因为 $|PF_1|>|PF_2|$，说明离\textbf{右焦点}更近，故 $P$ 在右支。
由比例：
\[
 ex_0+a=3(ex_0-a)
 \ \Longrightarrow\ 2ex_0=4a
 \ \Longrightarrow\ x_0=\frac{2a}{e}=\frac{32}{5}.
\]
代回曲线方程：
\[
 \frac{y_0^2}{9}
 =\frac{(32/5)^2}{16}-1
 =\frac{64}{25}-1=\frac{39}{25},
 \qquad y_0=\pm\frac{3\sqrt{39}}5.
\]
\[
 \boxed{P\left(\frac{32}{5},\ \pm\frac{3\sqrt{39}}5\right)}.
\]
\textbf{自检：} $ex_0=8$，两距离分别是 $12$ 和 $4$，比为 $3:1$，
差为 $8=2a$。左支不会满足题设的距离大小关系。

\medskip
\textbf{\textcolor{CExample}{例 4（新增：左支去绝对值，防止符号写反）}}

\textbf{题目：}
在双曲线 $\frac{x^2}{9}-\frac{y^2}{16}=1$ 上，
$P$ 的横坐标是 $x_0=-6$，求到左右焦点的距离。

\textbf{解：} $a=3$，$c=5$，$e=\frac53$。
$x_0=-6\le-3$，在左支；
$t=e|x_0|=\frac53\cdot6=10$。
左焦点\textbf{同侧短}、右焦点\textbf{异侧长}，因此
\[
 \boxed{|PF_1|=t-a=7},\qquad
 \boxed{|PF_2|=t+a=13}.
\]
\textbf{自检：} 两距离差为 $6=2a$，和为 $20=2e|x_0|$。

\medskip
\textbf{\textcolor{CExample}{例 5（新增：焦半径反推横坐标）}}

\textbf{题目：}
点 $P$ 在双曲线 $\frac{x^2}{16}-\frac{y^2}{9}=1$
\textbf{左支}，且 $|PF_1|=6$，求横坐标 $x_0$。

\textbf{解：} 左焦点同侧短，$|PF_1|=-ex_0-a$，
其中 $e=\frac54$、$a=4$。于是
\[
 6=-\frac54x_0-4
 \ \Longrightarrow\ \boxed{x_0=-8}.
\]
检查 $-8\le-4$，确实位于左支；
另一焦半径为 $|PF_2|=6+2a=14$。

\medskip
\textbf{\textcolor{CExample}{例 6（新增：两个端点的低成本检验）}}

\textbf{题目：}
若 $P$ 分别取双曲线的左、右顶点，写出两焦半径。

\textbf{解：}
\[
 \begin{array}{c|cc}
   P&|PF_1|&|PF_2|\\ \hline
   (a,0)&c+a&c-a\\
   (-a,0)&c-a&c+a
 \end{array}
\]
\textbf{意义：} 这是考场检查公式下标与符号是否写反的最快办法。

\end{examplebox}
